\section{What cropping did}
\label{sec:cropping}

This section answers the question the crop was built to answer. It reports the
effect on behaviour first, then the effect on where the policies look, since
the two need not agree and here they do not. Figure~\ref{fig:rim} carries the
second.

On recordings it never saw, cropping changes almost nothing. Over the shared
horizon the action-chunking family improves by about one per cent and the
diffusion family worsens by under half a per cent. Section~\ref{sec:threats}
gives the sample these rest on, and both effects are smaller than that sample
can resolve. The honest reading is that this crop makes no difference either
family's held-out error can detect, not that it makes a small positive or
negative one.

One effect is larger than the noise, and only the horizon view exposes it. For
the action-chunking family the sign of the crop's effect depends on where in
the chunk the error is measured: cropped is better over the first thirty-two
steps and worse over the full hundred. A single averaged error therefore
reports whichever the horizon happened to favour, and the same pair of runs
supports "cropping helps" or "cropping hurts" according to a choice the reader
never sees. This is the clearest argument in the report for publishing a curve
rather than a scalar.

\begin{figure}[t]
  \centering
  \includegraphics[width=\textwidth]{rim_attention.png}
  \caption{Gradient mass falling within a band around the edge of the tactile
  images, pooled over the four tactile cameras, expressed relative to what a
  flat map would put in the same band. The dashed line is that flat map. Values
  below it mean the policy puts \emph{less} weight at the edges than an
  indifferent model would. Three band widths are shown because the conclusion
  should not depend on where the band is drawn.}
  \label{fig:rim}
\end{figure}

Figure~\ref{fig:rim} tests the original hypothesis directly rather than through
task error, and it does not support it. Neither family puts more weight on the
rim than a flat map would: every uncropped bar sits between two thirds and four
fifths of the indifferent value. A model reading light leakage at the gel
boundary would sit above the line, not below it. The result holds at all three
band widths, so it is not an artefact of where the band was drawn.

Cropping moves both families \emph{toward} the rim rather than away from it,
which is the opposite of the intended effect and is not mysterious. Removing
the outer rows and stretching the remainder means the surviving pixels nearest
the edge were previously further in, so an unchanged attention pattern
mechanically lands closer to the new boundary. The measure cannot distinguish
that from a genuine change in what the policy seeks, and neither can these
runs.

Two limits on this figure. The two families produce gradient maps at different
spatial resolutions, coarse enough in one case that a narrow band cannot be
resolved at all, so comparing one family against the other is fair only at the
widest band; comparing a family against its own cropped twin is fair at every
band, and that is the comparison the figure is for. And the frames behind it
are drawn from recordings the policies trained on, because that is what the
attribution runs covered --- so it describes where a trained policy looks on
familiar material, and Section~\ref{sec:inputs} returns to what that can and
cannot support.

\paragraph{Separating the rim from the stretch.}
The crop removes the rim and stretches what remains, so every result above
is about both. A third arm of the action-chunking family separates them: it
crops exactly as the others do and does not resize, so its fingertip images
arrive a fifth shorter than its overhead image. It shares every other setting
with its two twins, including the random seed, the batch, the number of steps
and the split. This family is the only one that allows it. The diffusion
implementation refuses cameras of different shapes outright, and the
fine-tuned model would pad a shorter image rather than accept it, which is a
different change again.

\begin{table}[t]
  \centering
  \caption{Three arms of the action-chunking family that differ only in what
  is done to the fingertip images: nothing, the crop with the stretch back to
  full height, or the crop alone. Root-mean-square action error, over the
  first ten and thirty-two planned steps and over the whole hundred. Each arm
  is one training run.}
  \label{tab:stretch}
  \small
  \begin{tabular}{lrrrr}
\toprule
& \multicolumn{3}{c}{held-out} & trained-on \\
\cmidrule(lr){2-4}\cmidrule(lr){5-5}
ACT arm & first 10 & first 32 & all 100 & all 100 \\
\midrule
not cropped & 7.88 & 10.11 & 14.09 & 2.67 \\
cropped and stretched back & 7.88 & 9.99 & 14.68 & 2.63 \\
cropped, not stretched & 7.58 & 9.56 & 14.08 & 2.59 \\
\bottomrule
\end{tabular}

\end{table}

\textbf{Without the stretch, the crop is never worse than either other arm, and
is the best of the three over the near horizons.} Table~\ref{tab:stretch} shows it a few per cent lower than the
uncropped arm over the near horizons and level with it over the whole
hundred steps, while the stretched arm's penalty over the whole hundred ---
the one number in which cropping looked harmful --- disappears once the stretch
is removed. That pattern is what one would see if the stretch, not the loss of
the rim, were what cost the original crop at long range. It is a direction and
not a result: each arm is a single training run, and how far a second training
seed would move each row has not been measured. This family plans without
sampling, so there is no sampler floor to set against these differences either.
